\chapter{商业模式}

\section{商业模式整体逻辑}

本项目采取「\hlgreen{免费工具引流 + 精准服务变现}」的双轮驱动模式：以免费的拍照识别与预警服务作为获客入口，快速聚集小农户用户与数据；在此基础上，通过面向 B 端的农资渠道、品牌推广、政府采购与面向 C 端的增值会员实现多元化变现。整体逻辑为：\hlgold{用极简工具换用户，用用户规模换渠道价值，用数据壁垒换长期定价权}。

\section{收入来源设计}

\subsection{农资渠道佣金分成}

当系统识别并给出用药方案后，引导用户跳转合作农资电商或线下门店购买对应农资，按成交额抽取佣金分成。这是与产品价值链条最贴合、随用户规模线性增长的收入来源。

\subsection{农药品牌精准推广合作}

基于「病虫害—用药方案」的精准匹配场景，向农药、肥料品牌提供精准推广服务：在特定作物、特定病虫害、特定区域的防治方案中，合规展示合作品牌产品，按曝光或成交计费。该模式转化效率远高于泛流量广告。

\subsection{增值农技服务会员}

为有更高需求的规模经营户提供增值服务，包括：历史病害档案、多地块管理、专家在线问诊、个性化防治历、优先预警等，按年收取会员费。

\subsection{政府农业服务项目采购}

面向地方政府农业部门、供销社、植保站，提供区域病虫害监测预警平台、统防统治数字化工具与数据服务，以项目采购或年度服务费形式变现。

\section{成本结构分析}

本项目成本以\hlgreen{技术研发与模型迭代}和\hlgreen{市场推广与渠道建设}为主，辅以数据采集标注、人力管理与运营开支（详见第九章财务规划）。由于采用端侧推理，模型推理的边际算力成本趋近于零，规模效应显著。

\section{盈利路径与阶段规划}

\begin{enumerate}[leftmargin=1.6em]
  \item \hlgreen{第 1 年（市场验证期）}：以免费产品打透 1--2 个农业大县的样板市场，验证「识别—方案—预警」闭环与用户留存，同步启动农资佣金与政府项目试点，实现首批收入。本阶段\hlgold{不以营收为目标}，核心是沉淀田间数据与验证农户付费意愿。
  \item \hlgreen{第 2 年（规模扩张期）}：把样板县经验复制到周边县域、逐步铺开渠道，佣金与推广收入快速增长，会员与政府项目形成稳定现金流，于第 2 年末接近单月盈亏平衡。
  \item \hlgreen{第 3 年（商业化兑现期）}：数据壁垒与区域病害地图成型，四类收入稳步放量，实现整体盈利，并具备向更大区域复制与对外融资的条件。
\end{enumerate}

\noindent\textit{\small 说明：本项目三阶段规划刻意按「小步验证、稳健扩张」的节奏设计，各阶段收入目标均以学生团队实际可投入的人力与资源为约束条件测算，不追求短期的规模扩张。}
